Escribiendo \(M_i=V_ic_i\), los balances son
\[
M_1'=5+5c_2-15c_1-k_1M_1,\qquad
M_2'=10+5c_1-25c_2-k_2M_2.
\]
Con los datos iniciales de limpieza,
\[
c_1'=0.1-0.4c_1+0.1c_2,\qquad
c_2'=0.1+0.05c_1-0.25c_2.
\]
En la concentración inicial común, las pendientes son \(0.04\) y \(0.06\,\mathrm{mg/(m^3\,min)}\), por lo que la igualdad no persiste. Para que persista desde cualquier valor común se necesitan
\[
\frac{s_1}{V_1}=\frac{s_2}{V_2},\qquad
\frac{Q_1}{V_1}+k_1=\frac{Q_2}{V_2}+k_2.
\]
Aquí basta establecer \(k_2=0.10\,\mathrm{min}^{-1}\). La concentración común satisface entonces \(c'=0.1-0.3c\), con \(c(0)=0.20\).
